Procesní hodnocení klade důraz na zdokumentovaný postup, kroky a odůvodnění místo jen konečného výsledku; některé nástroje (např. „Pokrok v matematice“) podporují požadavek nahraného postupu \cite{kb_ai_101_tipu_pdf_04e4a115_page_8_1}.

Důsledky pro návrh zadání a rubriky:
\begin{itemize}
  \item Požadujte jasné mezikroky a zdůvodnění (sken/text/video); pokud možno považujte nahraný postup za povinný artefakt \cite{kb_ai_101_tipu_pdf_04e4a115_page_8_1}.
  \item V rubrice uveďte váhy pro kvalitu postupu, správnost metody a reflexi chyb (např. postup 50\%, výsledek 20\%, reflexe 30\%).
  \item Naplánujte milníky (inkrementální odevzdání) a spot‑checky pro formativní zpětnou vazbu.
\end{itemize}

Praktický příklad (matematika):
\begin{enumerate}\item Zadání: řešte modelový problém a nahrajte pracovní postup (fotografie rukopisu nebo export z OneNote) s komentáři k rozhodnutím.\item Hodnocení: nástroj označí pravděpodobné chyby; učitel provede spot‑check 5–10\,\% prací a finálně schválí bodové ohodnocení \cite{kb_ai_101_tipu_pdf_04e4a115_page_8_1}.\end{enumerate}

Rychlý checklist pro vyučující:
\begin{itemize}
  \item Je postup požadován jako samostatný artefakt? Obsahuje rubrika váhy pro procesní kritéria? Jsou naplánovány milníky a spot‑checky?
\end{itemize}